% from ega2-1-fr.tex
\oldpage[II]{48}

\begin{env}[2.8.15]
\label{II.2.8.15-fr}
Les notations et hypothèses étant celles de
(\hyperref[II.2.8.1-fr]{2.8.1}), il résulte de
(\hyperref[II.2.4.7-fr]{2.4.7}) que l'on ne change pas le morphisme $\Phi$,
à des isomorphismes près, en remplaçant $S_0$ et $S'_0$ par
$\mathbf{Z}$ et $\varphi_0$ par l'application identique, puis en remplaçant
$S$ et $S'$ par $S^{(d)}$ et ${S'}^{(d)}$ respectivement ($d>0$) et
$\varphi$ par sa restriction $\varphi^{(d)}$ à $S^{(d)}$.
\end{env}

\subsection{Sous-schémas fermés d'un schéma $\operatorname{Proj}(S)$.}
\label{subsection:II.2.9-fr}

\begin{env}[2.9.1]
\label{II.2.9.1-fr}
Si $\varphi:S\to S'$ est un homomorphisme d'anneaux gradués, on dit que
$\varphi$ est (TN)-\emph{surjectif} (resp. (TN)-\emph{injectif},
(TN)-\emph{bijectif}) s'il existe un entier $n$ tel que, pour $k\geq n$,
$\varphi_k:S_k\to S'_k$ soit surjectif (resp. injectif, bijectif). Alors il
résulte de (\hyperref[II.2.8.15-fr]{2.8.15}) que l'étude de $\Phi$ peut se
ramener au cas où $\varphi$ est \emph{surjectif} (resp. \emph{injectif},
\emph{bijectif}). Au lieu de dire que $\varphi$ est (TN)-bijectif, on dit aussi
que c'est alors un (TN)-\emph{isomorphisme}.
\end{env}

\begin{proposition}[2.9.2]
\label{II.2.9.2-fr}
Soit $S$ un anneau gradué à degrés positifs et soit
$X=\operatorname{Proj}(S)$.
\begin{enumerate}
  \item[\rm(i)] Si $\varphi:S\to S'$ est un homomorphisme (TN)-surjectif
    d'anneaux gradués, le morphisme correspondant $\Phi$
    (\hyperref[II.2.8.1-fr]{2.8.1}) est défini dans
    $\operatorname{Proj}(S')$ tout entier et est une immersion fermée de
    $\operatorname{Proj}(S')$ dans $X$. Si $\mathfrak{J}$ est le noyau de
    $\varphi$, le sous-schéma fermé de $X$ associé à $\Phi$ est défini par
    le faisceau d'idéaux quasi-cohérent $\widetilde{\mathfrak{J}}$ de
    $\mathcal{O}_X$.
  \item[\rm(ii)] Supposons en outre l'idéal $S_+$ engendré par un nombre
    fini d'éléments homogènes de degré~$1$. Soit $X'$ un sous-schéma fermé
    de $X$, défini par un faisceau quasi-cohérent d'idéaux $\mathcal{J}$ de
    $\mathcal{O}_X$. Soit $\mathfrak{J}$ l'idéal gradué de $S$, image
    réciproque de $\Gamma_\bullet(\mathcal{J})$ par l'homomorphisme canonique
    $\alpha:S\to\Gamma_\bullet(\mathcal{O}_X)$
    (\hyperref[II.2.6.2-fr]{2.6.2}), et posons
    $S'=S/\mathfrak{J}$. Alors $X'$ est le sous-schéma associé à
    l'immersion fermée $\operatorname{Proj}(S')\to X$ correspondant à
    l'homomorphisme canonique d'anneaux gradués $S\to S'$.
\end{enumerate}
\end{proposition}

\begin{proof}
\begin{enumerate}
  \item[\rm(i)] On peut supposer $\varphi$ surjectif
    (\hyperref[II.2.9.1-fr]{2.9.1}). Comme par hypothèse
    $\varphi(S_+)$ engendre $S'_+$, on a
    $G(\varphi)=\operatorname{Proj}(S')$. D'autre part, la seconde assertion
    se vérifie localement sur $X$~; soit donc $f$ un élément homogène de
    $S_+$ et posons $f'=\varphi(f)$. Comme $\varphi$ est un homomorphisme
    surjectif d'anneaux gradués, on vérifie aussitôt que
    $\varphi_{(f')}:S_{(f)}\to S'_{(f')}$ est surjectif et que son noyau est
    $\mathfrak{J}_{(f)}$, ce qui achève de prouver (i)
    (\hyperref[I.4.2.3-fr]{I, 4.2.3}).
  \item[\rm(ii)] En vertu de (i), on est ramené à vérifier que
    l'homomorphisme $\widetilde{j}:\widetilde{\mathfrak{J}}\to\mathcal{O}_X$
    déduit de l'injection canonique $j:\mathfrak{J}\to S$ est un
    isomorphisme de $\widetilde{\mathfrak{J}}$ sur $\mathcal{J}$, ce qui
    résulte de (\hyperref[II.2.7.11-fr]{2.7.11}).
\end{enumerate}
\end{proof}

On notera que $\mathfrak{J}$ est le \emph{plus grand} des idéaux gradués
$\mathfrak{J}'$ de $S$ tels que
$\widetilde{j}(\widetilde{\mathfrak{J}'})=\mathcal{J}$, car on vérifie
immédiatement en remontant aux définitions
(\hyperref[II.2.6.2-fr]{2.6.2}) que cette relation implique
$\alpha(\mathfrak{J}')\subset\Gamma_\bullet(\mathcal{J})$.

\begin{corollary}[2.9.3]
\label{II.2.9.3-fr}
On suppose vérifiées les hypothèses de
(\hyperref[II.2.9.2-fr]{2.9.2}, (i)) et en outre que l'idéal $S_+$ est
engendré par $S_1$~; alors $\Phi^*(\widetilde{S(n)})$ est canoniquement
isomorphe à $\widetilde{S'(n)}$ pour tout $n\in\mathbf{Z}$, et par suite
$\Phi^*(\mathcal{F}(n))$ à $\Phi^*(\mathcal{F})(n)$ pour tout
$\mathcal{O}_X$-Module $\mathcal{F}$.
\end{corollary}

\begin{proof}
C'est un cas particulier de (\hyperref[II.2.8.8-fr]{2.8.8}), compte tenu de
(\hyperref[II.2.5.10.2-fr]{2.5.10.2}).
\end{proof}

\begin{corollary}[2.9.4]
\label{II.2.9.4-fr}
On suppose vérifiées les hypothèses de
(\hyperref[II.2.9.2-fr]{2.9.2}, (ii)). Pour que le sous-préschéma fermé $X'$
de $X$ soit intègre, il faut et il suffit que l'idéal gradué
$\mathfrak{J}$ soit premier dans $S$.
\end{corollary}

